\chapter{المحيط والمساحة والحجم}\label{ch:g6:measure}

كم يبلغ طول السياج، وكم يبلغ اِتساع الحقل، وكم يسع الخزان من
الماء؟ ثلاثة أسئلة مختلفة، وثلاثة مقادير مختلفة ---
\omterm{def:g6:measure:perimeter}{المحيط} \omterm{def:g6:measure:area}{والمساحة} \omterm{def:g6:measure:volume}{والحجم} --- لكل منها وحداته. والخلط بينها
أشيع خطأ في الهندسة؛ وهذا الفصل يفصل بينها فصلًا
نهائيًّا.

\section{الأطوال والمحيط}

\begin{definition}[المحيط]\label{def:g6:measure:perimeter}
\emph{محيط}\index{محيط} الشكل هو الطول الكلي
لحدوده. ويُقاس بوحدات الطول: الميليمتر (مم)،
والسنتيمتر (سم)، والمتر (م)، والكيلومتر (كم)، مع
\[
1 \text{ كم} = 1000 \text{ م}, \qquad
1 \text{ م} = 100 \text{ سم}, \qquad
1 \text{ سم} = 10 \text{ مم}.
\]
\end{definition}

\begin{example}\label{ex:g6:measure:perimeter}
للمستطيل الذي طوله $L$ وعرضه $w$ \omterm{def:g6:measure:perimeter}{محيط}
\[
P = L + w + L + w = 2 \times (L + w).
\]
ومن أجل مستطيل $7$ سم في $4$ سم: $P = 2 \times (7 + 4) = 2 \times 11 =
22$ سم. وللمربع الذي طول ضلعه $c$ \omterm{def:g6:measure:perimeter}{محيط} $4c$.
\end{example}

\begin{proposition}[طول الدائرة]\label{prop:g6:measure:circle}
\omterm{def:g6:measure:perimeter}{محيط} الدائرة (أو \emph{طولها}) التي نصف قطرها $r$ هو
\[
P = 2 \pi r ,
\]
حيث $\pi \approx 3.14$ هو العدد نفسه من أجل كل دائرة.
\end{proposition}
\admitted

\begin{example}
بركة دائرية نصف قطرها $5$ م. فحدّها يقيس
$2 \times \pi \times 5 = 10\pi \approx 31.4$ م. واِحتفظ بالقيمة المضبوطة
$10\pi$ أطول ما يمكن؛ ولا تقرّب إلا في النهاية.
\end{example}

\section{المساحات}

\begin{definition}[المساحة]\label{def:g6:measure:area}
\emph{مساحة}\index{مساحة} الشكل تقيس السطح الذي تغطّيه:
أي كم مربع وحدة يدخل فيه. والوحدات: سم$^2$ (مربع طول ضلعه
$1$ سم)، وم$^2$، وكم$^2$، \dots\ واِنتبه:
\[
1 \text{ م}^2 = 100 \times 100 \text{ سم}^2 = 10\,000 \text{ سم}^2
\]
--- فالمتر المربع مربع $100$ سم في $100$ سم، ومنه كل درجة من
سلّم الوحدات تساوي $100$، لا $10$.
\end{definition}

\begin{omfigure}
\begin{tikzpicture}[scale=0.68]
  \draw[gray!50, thin] (0,0) grid (7,4);
  \draw[very thick, omDef] (0,0) rectangle (7,4);
  \fill[omThm!25] (0,3) rectangle (1,4);
  \draw[omThm, thick] (0,3) rectangle (1,4);
  \node[font=\small, omThm] at (2.6,3.5) {$1$ مربع وحدة};
  \node[font=\small] at (3.5,-0.5) {$7 \times 4 = 28$ مربع وحدة};
\end{tikzpicture}

{\small \omterm{def:g6:measure:area}{مساحة} المستطيل: $7$ أعمدة في كل منها $4$ مربعات وحدة،
إذن $7 \times 4 = 28$ مربعًا في المجموع. ولهذا كانت \omterm{def:g6:measure:area}{المساحة} $=$ الطول
$\times$ العرض.}
\end{omfigure}

\begin{proposition}[صيغ المساحات الأساسية]\label{prop:g6:measure:formulas}
\begin{center}
\renewcommand{\arraystretch}{1.5}
\begin{tabular}{l|l}
الشكل & \omterm{def:g6:measure:area}{المساحة} \\
\hline
مستطيل ($L \times w$) & $A = L \times w$ \\
مربع (طول الضلع $c$) & $A = c^2$ \\[4pt]
\omterm{def:g6:shapes:triangles}{مثلث قائم الزاوية} (ضلعا القائمة $a$ و $b$) & $A = \dfrac{a \times b}{2}$ \\[4pt]
قرص (\omterm{def:g3:shapes:shapes}{نصف القطر} $r$) & $A = \pi r^2$ \\
\end{tabular}
\end{center}
\end{proposition}

\begin{proof}[برهان المثلث القائم الزاوية]
نسختان من \omterm{def:g6:shapes:triangles}{مثلث قائم الزاوية} ضلعا قائمته $a$ و $b$، ملصوقتان على طول
الوتر، تكوّنان مستطيلًا $a \times b$. إذن \omterm{def:g6:measure:area}{مساحة} المثلث
نصف \omterm{def:g6:measure:area}{مساحة} المستطيل: أي $\frac{ab}{2}$. (وصيغة المستطيل هي
عدّ مربعات الوحدة أعلاه؛ وصيغة القرص مقبولة، اُنظر
\cref{ch:g7:areas}.)
\end{proof}

\begin{omfigure}
\begin{tikzpicture}[scale=0.75]
  \draw[thick, fill=omDef!15] (0,0) -- (4,0) -- (0,2.5) -- cycle;
  \draw[thick, omExo, dashed] (4,0) -- (4,2.5) -- (0,2.5);
  \draw[thin] (0.3,0) -- (0.3,0.3) -- (0,0.3);
  \node[below, font=\small] at (2,-0.1) {$a$};
  \node[left, font=\small] at (0,1.25) {$b$};
  \node[font=\small, omExo] at (2.8,1.7) {النسخة الثانية};
\end{tikzpicture}

{\small \omterm{def:g6:shapes:triangles}{المثلث القائم الزاوية} نصف مستطيل: ومن ثمّ الصيغة
$\frac{a\times b}{2}$.}
\end{omfigure}

\begin{example}[الأشكال المركّبة]\label{ex:g6:measure:composite}
غرفة على شكل حرف L هي مستطيل $6$ م $\times$ $4$ م حُذف منه ركن مربع $2$ م
$\times$ $2$ م. ومساحتها، خطوة خطوة:
\begin{enumerate}
  \item المستطيل الكامل: $6 \times 4 = 24$ م$^2$؛
  \item المربع المحذوف: $2 \times 2 = 4$ م$^2$؛
  \item \omterm{def:g6:measure:area}{المساحة} الباقية: $24 - 4 = 20$ م$^2$.
\end{enumerate}
وأما \emph{محيطها} فليس $20$ من أيّ شيء: فبالدوران حول حرف L،
ما زال الحدّ يقيس $6 + 4 + 6 + 4 = 20$ م --- فقَطعا الركن
يعوّضان \omterm{def:g6:lines:objects}{قطعتين} متساويتين من الجدار. العدد نفسه بالمصادفة،
لكنه أمتار مربعة في واحدة، وأمتار في الأخرى!
\end{example}

\begin{remark}[المحيط نفسه، ومساحات مختلفة]
يمكن أن يكون لشكلين \omterm{def:g6:measure:perimeter}{المحيط} نفسه ومساحتان مختلفتان جدًّا:
فالمربع $5 \times 5$ والمستطيل $9 \times 1$ كلاهما محيطه
$20$، لكن مساحتيهما $25$ و $9$. \omterm{def:g6:measure:perimeter}{فالمحيط} \omterm{def:g6:measure:area}{والمساحة} مقداران مستقلان
حقًّا.
\end{remark}

\section{الأحجام}

\begin{definition}[الحجم]\label{def:g6:measure:volume}
\emph{حجم}\index{حجم} \omterm{def:g5:solids:def}{المجسّم} يقيس الحيز الذي يملؤه:
أي كم \omterm{def:g5:solids:def}{مكعب} وحدة يدخل فيه. والوحدات: سم$^3$، وم$^3$، \dots\ مع
$1$ م$^3 = 1\,000\,000$ سم$^3$ (وكل درجة من السلّم تساوي
$1000$). وللسوائل تُستعمل أيضًا وحدة \emph{اللتر}:
\[
1 \text{ لتر} = 1 \text{ دسم}^3 = 1000 \text{ سم}^3,
\qquad
1 \text{ م}^3 = 1000 \text{ لتر}.
\]
\end{definition}

\begin{proposition}[حجم الصندوق]\label{prop:g6:measure:box}
الصندوق (أي \emph{متوازي المستطيلات}) الذي طوله $L$ وعرضه $w$
وارتفاعه $h$ حجمه
\[
V = L \times w \times h .
\]
\end{proposition}

\begin{proof}
تحوي الطبقة السفلى $L \times w$ \omterm{def:g5:solids:def}{مكعب} وحدة
(صورة \cref{def:g6:measure:area}، بالمكعبات)، وهناك $h$
طبقة كهذه.
\end{proof}

\begin{omfigure}
\begin{tikzpicture}[scale=0.6]
  \foreach \y in {0,1} {
    \foreach \x in {0,...,3} {
      \foreach \z in {0,1} {
        \begin{scope}[shift={({\x+0.45*\z},{\y+0.3*\z})}]
          \draw[fill=omDef!12] (0,0) rectangle (1,1);
          \draw[fill=omDef!20] (0,1) -- (0.45,1.3) -- (1.45,1.3) -- (1,1)
            -- cycle;
          \draw[fill=omDef!28] (1,0) -- (1.45,0.3) -- (1.45,1.3) -- (1,1)
            -- cycle;
        \end{scope}
      }
    }
  }
  \node[below, font=\small] at (2.5,-0.3)
    {$4 \times 2 \times 2 = 16$ مكعب وحدة};
\end{tikzpicture}

{\small صندوق مملوء بمكعبات وحدة: $4 \times 2$ مكعبات في الطبقة، وطبقتان
اثنتان.}
\end{omfigure}

\begin{example}\label{ex:g6:measure:aquarium}
حوض سمك قياساته $60$ سم في $30$ سم في $40$ سم (الارتفاع). \omterm{def:g6:measure:volume}{والحجم}:
\[
V = 60 \times 30 \times 40 = 72\,000 \text{ سم}^3 = 72 \text{ دسم}^3
= 72 \text{ لتر}.
\]
والتحويل خطوة خطوة: $1000$ سم$^3$ تصنع لترًا واحدًا،
ثم $72\,000 \div 1000 = 72$.
\end{example}

\section{تمارين}

\begin{exercise}[$\star$]\label{exo:g6:measure:1}
حوّل: $3.5$ م إلى سم؛ و $420$ مم إلى سم؛ و $0.8$ كم إلى م؛
و $25\,000$ م إلى كم.
\end{exercise}

\begin{exercise}[$\star$]\label{exo:g6:measure:2}
اِحسب \omterm{def:g6:measure:perimeter}{محيط}: مستطيل $8$ سم $\times$ $3.5$ سم؛
ومربع طول ضلعه $6.2$ سم؛ ومثلث أضلاعه $5$ سم و $7$ سم و
$9$ سم.
\end{exercise}

\begin{exercise}[$\star$]\label{exo:g6:measure:3}
مضمار جري دائري نصف قطره $50$ م. فكم يبلغ طول اللفّة الواحدة
(القيمة المضبوطة بالعدد $\pi$، ثم مقرّبة إلى المتر)؟ وكم لفّة تصنع
$2$ كم على الأقلّ؟
\end{exercise}

\begin{exercise}[$\star$]\label{exo:g6:measure:4}
اِحسب \omterm{def:g6:measure:area}{مساحة}: مستطيل $9$ سم $\times$ $4$ سم؛ ومربع طول
ضلعه $7$ م؛ \omterm{def:g6:shapes:triangles}{ومثلث قائم الزاوية} ضلعا قائمته $6$ سم و $10$ سم؛ وقرص
نصف قطره $3$ سم (القيمة المضبوطة، ثم مقرّبة إلى سم$^2$).
\end{exercise}

\begin{exercise}[$\star$]\label{exo:g6:measure:5}
حوّل: $3$ م$^2$ إلى سم$^2$؛ و $45\,000$ سم$^2$ إلى م$^2$؛ و $2.5$ لتر
إلى سم$^3$؛ و $4\,500$ لتر إلى م$^3$.
\end{exercise}

\begin{exercise}[$\star$]\label{exo:g6:measure:6}
حقل مستطيل طوله $120$ م وعرضه $85$ م. فكم مترًا من السياج
يلزم لإحاطته؟ وما مساحته؟
\end{exercise}

\begin{exercise}[$\star$]\label{exo:g6:measure:7}
اِحسب \omterm{def:g6:measure:volume}{حجم} صندوق $5$ سم $\times$ $4$ سم $\times$ $10$ سم،
\omterm{def:g6:measure:volume}{وحجم} \omterm{def:g5:solids:def}{مكعب} طول حرفه $3$ سم.
\end{exercise}

\begin{exercise}[$\star$]\label{exo:g6:measure:8}
اُرسم مستطيلين مختلفين \omterm{def:g6:measure:perimeter}{محيط} كل منهما $16$ سم، واِحسب
مساحتيهما. فأيّ مستطيليك مساحته أكبر؟
\end{exercise}

\begin{exercise}[$\star\star$]\label{exo:g6:measure:9}
شكل على هيئة حرف T مصنوع من مستطيل أفقي $10 \times 2$ فوق
مستطيل عمودي $2 \times 6$ (القياسات بالسنتيمتر). اِحسب
مساحته، ثم محيطه (دُر حول الحدّ بعناية، جامعًا
كل حرف).
\end{exercise}

\begin{exercise}[$\star\star$]\label{exo:g6:measure:10}
مسبح صندوق مستطيل طوله $25$ م، وعرضه $10$ م، وعمقه
$2$ م.
\begin{enumerate}
  \item فكم مترًا مكعبًا من الماء يسع حين يمتلئ؟
  \item وكم لترًا يساوي ذلك؟
  \item يُملأ المسبح بمعدّل $50\,000$ لتر في الساعة. فكم يستغرق
        الملء؟
\end{enumerate}
\end{exercise}

\begin{exercise}[$\star\star$]\label{exo:g6:measure:11}
حديقة مربعة طول ضلعها $20$ م فيها بركة دائرية نصف
قطرها $3$ م. فما \omterm{def:g6:measure:area}{مساحة} العشب فيها (القيمة المضبوطة، ثم مقرّبة
إلى م$^2$)؟
\end{exercise}

\begin{exercise}[$\star\star\star$]\label{exo:g6:measure:12}
لوح شوكولاتة قياساته $15$ سم $\times$ $6$ سم $\times$ $1$ سم.
ويضاعف الصانع الأبعاد \emph{الثلاثة} كلها.
\begin{enumerate}
  \item فبكم يُضرب \omterm{def:g6:measure:volume}{الحجم}؟ تحقّق بحساب
        الحجمين.
  \item والثمن يُضرب في $4$. فهل اللوح الكبير صفقة أفضل
        من الصغير؟
\end{enumerate}
\end{exercise}

\section{مسألة: سياج الملكة ديدو}

\begin{problem}[{مسألة نهاية الأسبوع --- بطول سياج ثابت،
أيّ شكل يحيط بأكبر أرض؟ المربع يغلب كل
مستطيل، والدائرة تغلب المربع، وجدار مخزن يقلب
كل شيء}]
\label{pb:g6:measure:1}
تقول الأسطورة إن الملكة ديدو، حين نزلت الساحل، مُنحت
«من الأرض ما يستطيع جلد ثور أن يحيط به» --- فقطّعت
الجلد شريطًا واحدًا رفيعًا طويلًا جدًّا وأحاطت بأرض
تكفي لتأسيس مدينتها. ومسألتها اليوم مسألتك: فلاح
يملك $20$ م من السياج بالضبط. وقد بيّن \cref{exo:g6:measure:8}
أن حظيرتين لهما \omterm{def:g6:measure:perimeter}{المحيط} نفسه قد تختلف
مساحتاهما؛ وتجد هذه المسألة الحظيرة \emph{الأفضل} --- وتكتشف
أن الجواب يتغيّر تمامًا حين يمنح جدار مخزن ضلعًا
مجّانًا.

\medskip
\noindent\textbf{الجزء الأول --- عشرون مترًا من السياج.}
\begin{enumerate}
  \item يجب أن تكون الحظيرة مستطيلًا يستعمل $20$ م من السياج كلها.
        تحقّق من أن حظيرة $1 \times 9$ وحظيرة $2 \times 8$ كلتيهما
        صالحتان، واِحسب مساحتيهما.
  \item اِشرح لماذا يكون \omterm{def:g1:addition:def}{مجموع} طول كل حظيرة صالحة وعرضها
        $10$ م. ثم اِصنع الجدول الكامل للحظائر ذات
        الأبعاد الصحيحة (من $1 \times 9$ إلى $5 \times 5$) مع
        مساحاتها. فأيّها الأفضل؟
  \item وهل تستحق الأضلاع العشرية التجربة؟ اِحسب مساحتَي
        الحظيرة $4.5 \times 5.5$ والحظيرة $4.9 \times 5.1$،
        وقارنهما بالمربع $5 \times 5$. فماذا
        تخمّن؟
  \item حوّل السؤال 2 مسألة السياج إلى سؤال عددي
        صرف: \emph{من بين كل أزواج الأعداد التي مجموعها
        $10$، أيّ زوج جداؤه الأكبر؟} أجب عنه من
        جدولك، واُذكر القاعدة العامّة التي تلاحظها.
  \item وهذا سبب كون الابتعاد عن المربع خسارة دائمًا،
        بالمقصّ لا بالجبر: اِنطلق من
        الحظيرة $5 \times 5$ وحوّلها إلى الحظيرة $6 \times 4$
        بنزع شريط ولصق آخر مكانه. فأيّ
        شريط يُنزع، وأيّهما يُضاف، ولماذا تخسر
        المبادلة مترًا مربعًا واحدًا بالضبط؟ واِشرح لماذا تخسر
        خطوة أخرى (إلى $7 \times 3$) أكثر من ذلك.
\end{enumerate}

\medskip
\noindent\textbf{الجزء الثاني --- المخزن، والدائرة.}
\begin{enumerate}[resume]
  \item تُبنى الحظيرة الآن ملاصقة لجدار مخزن طويل: فالجدار
        يعوّض أحد أطوال الحظيرة، وأمتار السياج $20$
        تغطّي الأضلاع الثلاثة الأخرى فقط. اِصنع جدول
        الحظائر ذات الأبعاد الصحيحة (العرض $w$، والطول $L$ على طول الجدار،
        $w + L + w = 20$) ومساحاتها. فأيّ حظيرة تفوز الآن ---
        وهل هي مربع؟
  \item اِشرح الفائزة بمرآة
        (\cref{ch:g6:symmetry}): اِعكس الحظيرة بالنسبة إلى جدار
        المخزن وتأمّل الحظيرة المضاعفة. فكم سياجًا
        تستعمل الحظيرة المضاعفة، وأيّ حظيرة مضاعفة هي الأفضل حسب الجزء الأول،
        وماذا يجعل ذلك الحظيرة الأصلية؟
  \item لم تبنِ ديدو مستطيلًا. فاثنِ أمتار السياج $20$
        دائرةً: وباستعمال $\pi \approx 3.14$، اِحسب
        نصف قطرها (إلى السنتيمتر)، ثم مساحتها (إلى م$^2$)
        (\cref{prop:g6:measure:circle}،
        \cref{prop:g6:measure:formulas}). وقارنها بالمربع
        $5 \times 5$.
  \item والمسألة العكسية: مطلوب حظيرة مساحتها $36$ م$^2$
        بالضبط، بأقلّ سياج ممكن. قارن
        المستطيلات $1 \times 36$ و $2 \times 18$ و $3 \times 12$
        و $4 \times 9$ و $6 \times 6$: فما المحيطات؟ وأيّ شكل
        هو الأرخص، وكيف يكون هذا السؤال صورة الجزء الأول
        في المرآة؟
  \item في جملة واحدة: لماذا تكون العلب والأنابيب وخزانات الماء
        \emph{مدوّرة} في كثير من الأحيان؟
\end{enumerate}

\medskip
\noindent\textbf{الجزء الثالث --- \omterm{def:g6:measure:perimeter}{المحيط} \omterm{def:g6:measure:area}{والمساحة} غريبان.}
\begin{enumerate}[resume]
  \item جِد مستطيلًا محيطه يتجاوز $100$ م لكن
        مساحته أصغر من $1$ م$^2$. (طويل جدًّا ورفيع
        جدًّا. \omterm{def:g6:decimals:places}{والأعداد العشرية} مسموحة.)
  \item صواب أم خطأ: «من بين شكلين، الذي محيطه
        أكبر مساحته أكبر.» أعطِ مثالًا مضادًّا
        من هذه المسألة.
  \item خذ المستطيل $6 \times 4$ واقطع فيه ثلمة مربعة
        $1 \times 1$ في وسط أحد ضلعيه الطويلين (والثلمة
        تفتح إلى الخارج، كسنّ مفقودة). اِحسب \omterm{def:g6:measure:area}{مساحة}
        الشكل المثلوم ومحيطه، وقارن الاثنين
        بالأصل. فماذا حدث لكل منهما؟
  \item يعرف رسّامو الخرائط واقعة غريبة: فكلما زادت تفاصيل الخريطة
        \emph{طال} قياس الساحل --- فكل تكبير
        يكشف ثلمات صغيرة جديدة، والسؤال 13 يبيّن ما
        تفعله كل ثلمة. اِشرح في جملة أو جملتين لماذا
        يكون «طول ساحل بلد ما» عددًا زلقًا،
        بينما «\omterm{def:g6:measure:area}{مساحة} ذلك البلد» ليست كذلك.
  \item اِمتحان الفلاح. مع $36$ م من السياج وجدار
        المخزن متاحًا، جِد أفضل حظيرة مستطيلة (اِستعمل
        مرآة السؤال 7)، وأعطِ مساحتها، وقارنها
        بأفضل حظيرة تُبنى بالسياج $36$ م في أرض مكشوفة. فكم
        يكسب جدار المخزن للفلاح؟
\end{enumerate}
\end{problem}
